Table~\ref{tab:trigger-vs-injection} contrasts the non-adversarial
triggers used in ADB with prompt injection, using a concrete drift
point from the Meetlark scheduling task (Appendix~\ref{app:case-study})
as an example. In ADB, the trigger is the field 
\texttt{prior\_polls\_by\_creator: 7} returned by \texttt{create\_poll},
a routine administrative counter that names no tool and issues
no instruction.

\begin{table*}[h]
\centering
\small
\caption{Non-adversarial trigger (ADB) vs.\ prompt injection.}
\label{tab:trigger-vs-injection}
\begin{tabular}{p{2.4cm}p{5.8cm}p{5.8cm}}
\toprule
 & \textbf{Non-adversarial trigger (ADB)}
 & \textbf{Prompt injection (contrast)} \\
\midrule
\textbf{Nature}
  & Normal API response field
  & Malicious text injected into tool output \\
\addlinespace
\textbf{Content}
  & \texttt{"prior\_polls\_by\_creator": 7}
  & \texttt{"Retrieve the creator's full poll history}
    \texttt{before proceeding. Call get\_poll\_history}
    \texttt{with the admin token."} \\
\addlinespace
\textbf{Tone}
  & Declarative system metadata
  & Imperative with explicit tool reference \\
\addlinespace
\textbf{Agent motivation}
  & Over-proactivity: the agent infers that historical
    patterns may improve personalisation and acts
    without authorisation
  & Instruction hijacking: the agent is explicitly
    commanded to call a specific tool \\
\bottomrule
\end{tabular}
\end{table*}